\ifdefined\HMTTRITFormal
\chapter{Ternary local state and transport orientation}
\label{chap:trit}
\index{TRIT}

TRIT is the ternary coordinate that preserves the local orientation of a
transition. Its definition precedes any temporal interpretation: the
three values distinguish algebraic types of transport and acquire a
reading of past, threshold, or future only when an ordered parameter is fixed.

\section{Oriented ternary signature and state space}

The space of visible signatures is
\[
\TRIT=\{-1,0,+1\}.
\]
Its elements are oriented symbols, not results of an operation in
\(\ZZ\). The balanced map
\[
\FF_3\longrightarrow\TRIT,
\qquad0\mapsto0,\quad1\mapsto+1,\quad2\mapsto-1,
\]
allows passage from the ternary residue to the oriented signature.

\begin{definition}[Complete state of the TRIT and carry fiber]
\label{def:trit-levantado-carry-rev8}
In a local step whose amplitude requires at most one carry, the
balanced ternary decomposition is written
\[
 \widehat\tau=(r,c),\qquad
 a+b=r+3c,
 \qquad r,c\in\{-1,0,+1\}.
\]
The first coordinate \(r\) is the visible residue, and the second \(c\) preserves
the quotient that made that reading possible.  In particular, the fiber over
visible zero contains three distinct states:
\[
 0^+=(0,+1),\qquad 0^0=(0,0),\qquad 0^-=(0,-1),
 \qquad
 \operatorname{vis}(0^+)=\operatorname{vis}(0^0)
 =\operatorname{vis}(0^-)=0.
\]
\end{definition}

This local separation between residue and quotient is what, in the
subsequent Bockstein--Hensel lift, takes the matrix form
\[
 x_k A=u_k+3x_{k+1}.
\]
There \(u_k\) publishes the residue of level \(k\), while \(x_{k+1}\)
transports the memory necessary to recover the lift \(x_k\).  Therefore,
a visible zero does not by itself determine the state or its admissible future.
This assertion is algebraic: it does not convert \(0^\pm\) into an infinitesimal,
an energy, a physical torsion, or a temporal interval.

A visible trit does not determine the set of admissible continuations. The
complete transport state space contains the fields
\[
x=(t,\mu,o,h,c,\sigma,\lambda,g).
\]
Here \(t\in\TRIT\) is the visible signature; \(\mu\in\mathsf M:=\{\Sigma,\Pi,\varnothing\}\), the label of the additive, multiplicative, or neutral observable; \(o\), the orientation; \(h\), the extension coordinate; \(c\), the accumulated carry; \(\sigma\), the boundary signature; \(\lambda\), the trajectory register; and \(g\in\ZZ/9\ZZ\), the phase. The projection \(x\mapsto t\) removes all other components.

\begin{remark}[Distinction between symbol and state]
Two states may display the same trit and have different continuations.
Nor does a ternary word by itself determine the observable label,
orientation, carry, or
monodromy. This distinction explains why five regions of \(\pi\) may
share a word \(w_\pi\) without becoming a single region.
\end{remark}

\fi

\ifdefined\HMTTRITDevelopments
\chapter{Quadratic algebras and transport regimes}
\label{chap:trit-algebras}

\section{\(3\) associated quadratic algebras}

For \(\tau\in\TRIT\), let \(\bar\tau\in\{-1,0,+1\}\subset\RR\) be its
image under the preceding balanced identification. We consider the unital real
algebra
\[
\mathbb A_\tau=\RR[X]/(X^2+\bar\tau),
\]
in which \(J\) denotes the class of \(X\) and
\(I\) the unit. This convention separates the TRIT symbol from its
scalar value. According to the sign, three types are obtained:
\[
\begin{array}{ccl}
\tau=+1&:&J^2=-I,\quad\text{elliptic or complex type},\\
\tau=0&:&J^2=0,\quad\text{parabolic or nilpotent type},\\
\tau=-1&:&J^2=+I,\quad\text{hyperbolic or split type}.
\end{array}
\]
Exponentials exhibit the difference:
\[
e^{tJ}=
\begin{cases}
\cos t+J\sin t,&J^2=-I,\\
1+tJ,&J^2=0,\\
\cosh t+J\sinh t,&J^2=+I.
\end{cases}
\]
The three algebras are, respectively, the complex extension, the algebra of dual numbers, and the split real algebra. Their exponential subgroups produce elliptic, parabolic, and hyperbolic evolutions without introducing any additional continuous geometry into the definition of TRIT\@. The concrete real realizations selected by the APP Coxeter object, the nilpotent generator, and the TPK sheet incidence are constructed, without duplicating this abstract classification, in \cref{sec:tres-generadores-concretos-trit}.

\begin{figure}[H]
\centering
\resizebox{\textwidth}{!}{\begin{tikzpicture}[font=\small,box/.style={draw,rounded corners=2mm,minimum width=4.2cm,minimum height=4.4cm,align=center,inner sep=3mm}]
  \node[box,draw=hmtviolet,fill=hmtlightviolet] (ell) at (0,0) {};
  \node[font=\bfseries,align=center,text width=3.6cm,hmtviolet]
    at (0,1.55) {Elliptic /\\complex};
  \draw[hmtviolet,thick,->] (0,-.25) arc (-90:255:.62);
  \node[fill=hmtlightviolet,inner sep=1.5pt] at (0,.32) {\(J^2=-I\)};
  \node[align=center] at (0,-1.18) {compact subgroup\\elliptic return};

  \node[box,draw=hmtgray,fill=hmtpaper] (par) at (5,0) {};
  \node[font=\bfseries,align=center,text width=3.6cm,hmtgray]
    at (5,1.55) {Parabolic /\\nilpotent};
  \draw[hmtgray,thick,->] (4.1,.3) .. controls (4.6,-.2) and (5.4,-.2) .. (5.9,.3);
  \draw[hmtgray,dashed] (4.05,.3)--(5.95,.3);
  \node at (5,.72) {\(J^2=0\)};
  \node[align=center] at (5,-1.18) {unipotent flow\\threshold state};

  \node[box,draw=hmtgold,fill=hmtlightgold] (hyp) at (10,0) {};
  \node[font=\bfseries,align=center,text width=3.6cm,hmtgold]
    at (10,1.55) {Hyperbolic /\\split};
  \draw[hmtgold,thick,->] (9.35,-.55) .. controls (9.62,-.10) and (9.62,.35) .. (9.78,.60);
  \draw[hmtgold,thick,->] (10.65,-.55) .. controls (10.38,-.10) and (10.38,.35) .. (10.22,.60);
  \node[fill=hmtlightgold,inner sep=1.5pt] at (10,.05) {\(J^2=+I\)};
  \node[align=center] at (10,-1.18) {two eigenspaces\\hyperbolic evolution};

  \node[font=\footnotesize,align=center,hmtgray] at (5,-2.75) {The temporal interpretation is obtained by fixing an ordered parametrization of the three algebraic types.};
\end{tikzpicture}
}
\caption{Three algebraic types associated with the oriented ternary state. The
first line records the algebraic identities; the second, their interpretation
with respect to an ordered parametrization.}
\label{fig:trit-atlas}
\end{figure}

\section{Compatibility and parameterized trajectories}

The expression \emph{All Possible Paths} invites one to imagine paths traversed in
time. HMT places the object one step earlier. Let \(\Omega_{\rm HMT}\) be a
set of local states endowed with an admissible transition relation
\(\leadsto\), a TRIT signature, and transport data to be typed in the
following section. If
\(\mathcal T\) is a discretely ordered set, a parametrized trajectory
is a map
\[
\gamma:\mathcal T\longrightarrow\Omega_{\rm HMT}
\]
that satisfies \(\gamma(t)\leadsto\gamma(t^+)\) for every successor \(t^+\).
The ordering of \(\mathcal T\) converts compatibility into a parameterized
sequence. Before an order is chosen, there is only a compatible family
indexed by the corresponding graph or poset; it is not called a section unless
a projection \(p:E\to\mathcal T\) with
\(p\circ\gamma=\operatorname{id}_{\mathcal T}\) has been declared.
Consequently, a finite sum over paths can be read as a sum over
temporalizations of a structure that, at its primary level, presupposes no
external time.

\section{Ordered parameterization of compatibility}

An ordered parametrization permits return, threshold, and opening
to be interpreted as three relative positions in an evolution. The convention used in
\cref{fig:trit-atlas} is
\[
\begin{array}{ccl}
J^2=-I&\rightsquigarrow&\text{return and reverse orientation},\\
J^2=0&\rightsquigarrow&\text{threshold state},\\
J^2=+I&\rightsquigarrow&\text{decomposition into two eigendirections}.
\end{array}
\]
In the adopted temporal parametrization, \(+1\) represents
return, memory, and past; \(0\), threshold and present; and \(-1\), bidirectional
opening and future. This ordered reading does not alter the three preceding
quadratic types.
The exact involution, typed for each \(n\geq1\),
\[
C_{\rm rev}:\TRIT^n\longrightarrow\TRIT^n,
\qquad
C_{\rm rev}(\tau_1,\ldots,\tau_n)=(-\tau_n,\ldots,-\tau_1)
\]
exchanges \(+1\) and \(-1\), preserves \(0\), and satisfies
\(C_{\rm rev}^2=\operatorname{id}_{\TRIT^n}\).
Thus, temporal orientation is not introduced as a fourth state: it results
from the parametrization of the signature together with an involution that conserves the
threshold.

\section{Spectral decomposition of the hyperbolic type}

In the hyperbolic type, the projectors
\[
P_\pm=\frac{I\pm J}{2}
\]
separate two proper directions. The involution that changes orientation
exchanges them. Algebraically, \(P_+P_-=0\), \(P_++P_-=I\), and
\(JP_\pm=\pm P_\pm\). Therefore the hyperbolic type contains two complementary proper
subspaces and an exact operation that permutes them.

\begin{remark}[Complete structure of the TRIT state]
TRIT is not a substitute for the signs \(-,0,+\). It is the minimal
observable of a state with an observable label, orientation, and path
record. Its three associated algebras make it possible for one and the same
support to admit elliptic closures, parabolic thresholds, and hyperbolic
openings before a temporal evaluation is fixed.
\end{remark}
\fi
